\section{Empirical Design}
\label{sec:empirical-design}

\subsection{Frozen pilot and reference labels}

The pilot comprises 100 constructed C++17 slow/fast pairs across 41 problem IDs and nine generator mechanisms. The slow programs include 38 state or work corruptions (A1), 24 omitted guards (A2), and 38 cases motivated by an algorithm replacement (B). A mechanism-informed AST diff sees both versions and assigns one reference macro-node in the slow program. The profiling, slicing, and ranking methods receive the slow program and its workload or profile. Reference agreement measures localization of the site designated by the benchmark. A fresh uninstrumented replay reproduced byte-identical stdout in 236 paired comparisons, covering all pairs at $N=16$ and $N=24$ and six additional sizes for the six loop-bound and stride cases.

We compare hit-only and work-weighted symptom rules on the same slow programs, then compare candidate sets from three slice configurations. An alert is an emitted symptom. Recall is the fraction of programs whose candidate set contains the reference. Reduction is the per-program quantity $1-|C|/|V|$, averaged over the 100 programs, where $C$ is the emitted set and $V$ is the canonical macro-node universe. Empty outputs contribute zero to recall and full reduction. The detector comparison uses only the constructed slow programs, so its alert rates measure sensitivity within the pilot.

\subsection{Ranking and dependence across examples}

We compare a symbolic role-priority policy refined on the pilot with a QLoRA-adapted Qwen2.5-Coder-1.5B ranker~\cite{dettmers2023qlora,hui2024qwen}. Both rank the same pruned slice. In the reference \nolinkurl{abstract_graph} setting, the neural prompt supplies AST roles, def/use cardinalities, hit and work exponents, and binary slice membership; Section~\ref{sec:method} defines the six comparison representations. The training objective predicts the single reference-node token. The frozen evaluation holds out groups defined by problem ID, identifier-normalized workload AST, and identical prompt tokens in any evaluated representation. Its 23 connected groups form five folds of 20 pairs. The grouped refit retains the pairs, references, and saved slices. It adds exact normalized source, identifier- and numeric-literal-normalized fingerprints on both versions, and exact binary Abstract prompt tokens. This yields 13 connected groups over five folds of 19, 19, 19, 21, and 22 pairs. The refit also evaluates the same quantized model and prompt without the adapter.

The refit uses double-quantized 4-bit NF4 weights and bfloat16 computation. Low-rank adapters in decoder blocks 20--27 use rank 16, LoRA $\alpha=32$, and dropout 0.05. Each fold trains for six epochs, with AdamW at learning rate $6\times10^{-4}$ and weight decay 0.01, microbatches of two, and four accumulation steps. The refit resets seed 42 per fold and recomputes forward activations. The historical evaluation used fold-dependent seeds and cached the first 20 blocks. Changes across those runs therefore reflect grouping, seed policy, and implementation together.

Top-1 counts an exact reference-node match. We report averages across pairs and averages across groups, weighting each represented group equally after averaging its constituent pairs. We also report results on the 66 slices with at least two candidates, since the 34 singleton slices contain the reference and force agreement in constrained ranking. Secondary paired tests flip the signs of whole-group differences. Descriptive percentile intervals resample groups in sorted ID order with seed \GroupedBootstrapSeed{} and sort \GroupedBootstrapDraws{} draws before taking nearest ranks. Leave-one-group-out endpoints are ranges, not intervals. These intervals describe sensitivity to group composition with fixed predictions. Cross-validation fits share training data across folds. The replay-only receipt \nolinkurl{development/grouped_statistics.json} gives denominators and Clopper--Pearson binomial reference intervals for pair rates.

\section{Empirical Results}
\label{sec:empirical-results}

\subsection{Work weighting and candidate selection}

Work weighting increased symptom coverage from 80 to 100 programs. The 20 additional alerts comprise nine pass-by-value container copies and eleven linear standard-library scans in loops. The work symptom matched 82 references, including all 62 in A2 and B. The pruned slice recovered all 18 remaining references and retained a mean of 1.90 nodes per program. The unpruned slice achieved the same recall with 4.89 candidates. The inter-scope rule alone missed seven references. Table~\ref{tab:eapg-pilot-slicing} reports these tradeoffs under the frozen profiling and slicing implementation.

\begin{table}[t]
\centering\small
\caption{Frozen candidate selection on 100 constructed pairs (38 A1, 24 A2, 38 B). Primary metric: reference retention (recall). Recall and mean reduction are percentages.}
\label{tab:eapg-pilot-slicing}
\begin{tabular}{@{}lrrr@{}}
\toprule
Candidate set & Recall & Mean size & Reduction \\
\midrule
Hit symptom & 62 & 0.80 & 95.2 \\
Work symptom & 82 & 1.00 & 94.1 \\
Inter-scope slice & 93 & 1.83 & 89.3 \\
Unpruned two-tier slice & 100 & 4.89 & 71.2 \\
Pruned two-tier slice & 100 & 1.90 & 89.0 \\
\bottomrule
\end{tabular}
\end{table}

\subsection{Representation, slicing, and ranking on the grouped split}

Under the 23-group frozen split, the Phase~1 slice-constrained \nolinkurl{abstract_graph} QLoRA model (\nolinkurl{6b-slice}) ranked two more references first than the pilot-refined symbolic policy by pair ($98/100$ vs.\ $96/100$) yet trailed by $2.17$ percentage points by group (95\% interval $\GroupedFrozenAllCI$; multi-candidate $64/66$ vs.\ $62/66$, $-3.33$ points $\GroupedFrozenMultiCI$). Under the stricter 13-group split removing shared template and prompt links (reusing the identically partitioned 5-fold \nolinkurl{pilot-resplit} checkpoints for \nolinkurl{abstract_graph} while training the remaining five representations under the same recipe), Table~\ref{tab:eapg-resplit-ablation} evaluates all six prompt serializations alongside the four fixed-slice non-neural baselines.

Three empirical findings emerge. First, without adaptation or slice constraints, base \texttt{Qwen2.5-Coder-1.5B} selects $0/66$ multi-candidate references across all six prompts ($1.71\%$--$14.66\%$ macro overall, with exact Top-1 logit ties on 2--45/100 pairs resolved by node order); applying slice constraints brings unadapted predictions ($43.41\%$--$58.83\%$) near uniform random guessing ($63.57/100$, $52.02\%$). Second, dynamic work exponents improve the pilot point estimate, but the grouped test does not distinguish the prompts: holding AST roles, D/U counts, hit exponents, and \texttt{SLICE} tags fixed between \nolinkurl{abstract_hits_only} and \nolinkurl{abstract_graph}, adding $\alpha_v^{\mathrm{work}}$ raises adapted Top-1 from $73/100$ ($60.09\%$) to $85/100$ ($77.49\%$) unconstrained and from $78/100$ ($72.32\%$) to $90/100$ ($84.44\%$) constrained (multi-candidate $44/66 \to 56/66$, $+15.75$ group points, $\GroupedHitsOnlyCI$), resolving hidden-work scans and algorithmic replacements in \texttt{A1.3}, \texttt{A1.4}, \texttt{A2.2}, and \texttt{B.2}.

Third, adding raw identifiers and code snippets lowers the pilot point estimate in a 1.5B adapter: \nolinkurl{full_eapg} drops to $75/100$ ($58.27\%$ macro) constrained and $41/66$ ($45.75\%$) multi-candidate ($0/6$ on \texttt{A1.4}, $7/13$ on \texttt{B.2}, $3/18$ off-symptom), trailing \nolinkurl{abstract_graph} by $26.17$ group points overall ($\GroupedFullAllCI$) and $34.02$ points on multi-candidate slices ($\GroupedFullMultiCI$), while \nolinkurl{static_eapg} drops to $39/100$ ($38.07\%$) unconstrained without dynamic slopes. Note also that identifier stripping collapses 80 distinct slow templates to 38 unique \nolinkurl{abstract_graph} prompts without creating cross-fold leakage, and \nolinkurl{abstract_graph} confuses stripped intra-loop resets in \texttt{A1.1} ($5/12$ constrained; $11/18$ off-symptom), where \nolinkurl{hotspot_raw} reads reset idioms directly and reaches $15/18$ ($75.00\%$ group) off-symptom. Finally, removing pilot loop-alias tuning (\nolinkurl{def_use=False}, matching prospective \nolinkurl{work_scoped}) lowers symbolic accuracy by 26 pairs across \texttt{B.1} and \texttt{B.2} to $70/100$ ($76.92\%$ group; $36/66$, $70.00\%$ multi-candidate). Under \nolinkurl{def_use=True}, the refined heuristic stays ahead of \nolinkurl{abstract_graph} ($96/100$ vs.\ $90/100$; $-10.23$ multi-candidate points, $\GroupedRefinedCI$), whereas both \nolinkurl{abstract_graph} and \nolinkurl{hotspot_raw} exceed unbiased \nolinkurl{def_use=False} overall while trailing on off-symptom group rate ($70.83\%$ and $75.00\%$ vs.\ $83.33\%$ for \nolinkurl{def_use=False}, despite \nolinkurl{hotspot_raw} hitting $15/18$ raw pairs vs.\ $14/18$).

\begin{table*}[t]
\centering\footnotesize
\setlength{\tabcolsep}{3.5pt}
\caption{Strict 13-group pilot ablation. Primary metric within this secondary analysis: equal-group Top-1 on multi-candidate slices. All neural models use Qwen2.5-Coder-1.5B (4-bit NF4).}
\label{tab:eapg-resplit-ablation}
\begin{tabular}{@{}lrrrrrrr@{}}
\toprule
Method / Representation & Tok & Base Unc.\ (All) & Base Con.\ (All) & QLoRA Unc.\ (All) & Constrained (All) & Multi-cand.\ (66) & Off-symp.\ (18) \\
\midrule
\texttt{raw\_code} & 246 & 9/100 (6.05) & 55/100 (48.78) & 54/100 (47.52) & 66/100 (67.84) & 32/66 (58.19) & 11/18 (67.50) \\
\texttt{hotspot\_raw} & 346 & 6/100 (4.27) & 52/100 (43.41) & 62/100 (50.35) & 83/100 (79.85) & 49/66 (73.81) & 15/18 (75.00) \\
\texttt{static\_eapg} & 413 & 3/100 (2.56) & 64/100 (46.51) & 39/100 (38.07) & 74/100 (66.43) & 40/66 (56.36) & 9/18 (50.00) \\
\texttt{abstract\_hits\_only} & 257 & 16/100 (11.18) & 66/100 (58.83) & 73/100 (60.09) & 78/100 (72.32) & 44/66 (64.02) & 9/18 (58.33) \\
\texttt{abstract\_graph} & 307 & 21/100 (14.66) & 62/100 (52.29) & 85/100 (77.49) & 90/100 (84.44) & 56/66 (79.77) & 11/18 (70.83) \\
\texttt{full\_eapg} & 496 & 2/100 (1.71) & 66/100 (47.91) & 73/100 (56.03) & 75/100 (58.27) & 41/66 (45.75) & 3/18 (23.33) \\
\midrule
Uniform guessing & -- & -- & -- & -- & 63.57/100 (52.02) & 29.57/66 (37.62) & 5.57/18 (29.38) \\
Raw slice order & -- & -- & -- & -- & 41/100 (46.15) & 7/66 (30.00) & 7/18 (50.00) \\
Symbolic (\texttt{def\_use=False}) & -- & -- & -- & -- & 70/100 (76.92) & 36/66 (70.00) & 14/18 (83.33) \\
Symbolic (\texttt{def\_use=True}) & -- & -- & -- & -- & 96/100 (92.31) & 62/66 (90.00) & 14/18 (83.33) \\
\bottomrule
\end{tabular}
\end{table*}

\subsection{Domain and scalar intake feasibility}

\paragraph{Raw-AST blockers on existing corpora}
All 100 archived slow programs fail the conventional scalar comparator's raw-AST admission (\texttt{0/100}): all 100 fail first on unsupported literal \texttt{\#include} headers, and after omitting top-level literal includes, all 100 contain multiple top-level functions (66 two-function, 34 three-function), with 87 failing first on multiple function definitions and 13 failing earlier on unsupported top-level declarations. On the pinned UOJ-Bench C++-family repair release~\cite{xu2026uoj} (714 pairs across 303 problems), only 11 pairs across 10 problems avoid unsupported literal includes (0/41 explicit C++17 pairs), and raw-AST CFG checks pass \texttt{0/714} both before and after hiding top-level literal includes; leading residual blockers are \texttt{using\_declaration} (281), raw parse errors (207), and \texttt{preproc\_def} (137).

\paragraph{Development-only 10-group OEIS scalar intake funnel}
To quantify scalar-domain attrition before recruiting independent authors, we evaluated trusted C++17 sketches across 10 tentative OEIS candidate groups (11 specs)~\cite{oeis2026} through GCC~15.1.0 syntax/preprocessing, the pinned prospective normalizer (\texttt{2c95df0e}), CFG reachability, reversible intervention mapping, finite-range growth bands, and closed-form operation counts. Four specs rejected prior to scaling because they require multi-statement Track~B rewrites (\texttt{A000788}, \texttt{A000045}, \texttt{A006257}) or alter output semantics (\texttt{A003586}); all 7 remaining specs passed syntax, CFG, and token mapping (\texttt{0} domain rejections). Two failed the finite-range slope floor ($\ge 0.55$) despite exact theoretical $0.50$-order speedups: Dirichlet quotient blocking \texttt{A006218\_quotient} ($O(\sqrt{N})\to O(N)$, slope $0.4827$) and trial-division totients \texttt{A002088} ($O(N^{1.5})\to O(N^2)$, slope $0.5116$). Five specs across five groups and four distinct mechanisms passed (\texttt{A001399}, \texttt{A000328}, \texttt{A006218\_step}, \texttt{A000201}, \texttt{A400004}); notably, \texttt{A000328} and \texttt{A000201} collapsed onto the same \texttt{monotone\_cursor\_reset} mechanism despite distinct normalized templates. On analytic profiles over the five eligible cases, \texttt{work\_scoped} retained 4/5 and hit 3/5 Top-1 (missing two-scope retention on \texttt{A400004} at $N7\to N9$, slope $0.6836$, and ranking sibling \texttt{N6} ahead of \texttt{N7} in AST order on \texttt{A001399}), while \texttt{work\_conventional} retained 5/5 and hit 0/5 Top-1 because CFG closure pulls zero-growth initializers and accumulators that the shared ranker hoists ahead of the loop symptom.

\subsection{Counterexamples, timing, and external contract diagnostics}

\begin{table}[t]
\centering\footnotesize
\setlength{\tabcolsep}{3pt}
\caption{Summary of domain audits, boundary diagnostics, and contract verification checks. Primary metric: number of cases meeting each check's stated criterion.}
\label{tab:dev-feasibility}
\begin{tabular}{@{}>{\raggedright\arraybackslash}p{0.28\columnwidth}>{\raggedright\arraybackslash}p{0.22\columnwidth}>{\raggedright\arraybackslash}p{0.44\columnwidth}@{}}
\toprule
Check / diagnostic & Cohort & Key quantitative outcome \\
\midrule
UOJ raw-AST audit & 714 pairs (303 prob.) & 0/714 pass CFG (11 header-clean) \\
OEIS scalar intake & 10 groups (11 specs) & 5 eligible (4 mech.); scoped 3/5, original CFG 0/5, nearest-preceding CFG \GroupedNearestCFG{} \\
PIE pair \texttt{p03497} & 18 runs ($6\times 3$) & Work $\alpha=2.00$ singleton; time slope $<0.50$ \\
Cursor reset & 1 family $\times$ 3 seeds & Shared \texttt{N5}; scoped 1/2 vs.\ conv.\ 2/5 \\
Comparator swap & 54 counting checks & Identical alerts on slopes $0.870$ vs.\ $-0.039$ \\
Transitive edit & 162 runs (9 builds) & Slope $0.899$; conv.\ 2/7 retains, scoped drops \\
Scope grid ($d\times k$) & 12 analytic cells & Scoped retains iff $d=0$; sizes $k+2$ vs.\ $d+6$ \\
Custody / contract & Mock 322; 7 statuses & 966 rows scored; all abstentions verified \\
\bottomrule
\end{tabular}
\end{table}

\paragraph{Human-written PIE validation pair}
Previously inspected development pair \texttt{p03497}~\cite{shypula2024pie} exposed an archived selector bug: an independent inner-loop formula matched counted traversals across all 18 slow inputs with $\alpha^{\mathrm{work}}=2.00$ across three seeds, yet the frozen rule selected the outer loop. The repaired selector chose the inner loop as a singleton out of 20 canonical nodes on every seed. Separate uninstrumented timing in \nolinkurl{runtime_validation_timing.json} yielded largest-size slow/fast median ratios of 11.21, 13.77, and 12.74 and log-ratio slopes of 0.398, 0.452, and 0.439 (a repeat check in \nolinkurl{RESEARCH_RECORD.md} gave 0.470--0.491), all below the 0.50 runtime confirmation threshold.

\paragraph{Executed cursor reset, comparator swap, and transitive edit}
On an authored cursor-reset fixture executed in Ubuntu WSL2 Bubblewrap, \texttt{work\_scoped} and \texttt{work\_conventional} shared symptom \texttt{N5} and retained edit \texttt{N4}, ranking it first of two vs.\ second of five across three seeds; keeping the slow program fixed while swapping a linear vs.\ quadratic comparator flipped the verified growth slope from $0.870$ ($37.267\times$) to $-0.039$ ($1.008\times$) without changing any single-side prediction. Conversely, on an authored transitive fixture passing the edited value through two intervening copies ($162$ plain/UBSan/counting runs across nine builds, slope $0.899$, ratio $47.200$), \texttt{work\_conventional} retained \texttt{N7} at rank two of seven behind shared symptom \texttt{N10} while \texttt{work\_scoped} emitted two nodes and dropped \texttt{N7}. Crossing copy depth $d\in\{0,1,2,4\}$ against overwritten writes $k\in\{0,2,8\}$ over six analytic sizes generalized both failures cleanly: scoped slicing retained the edit iff $d=0$ and emitted $k+2$ nodes, whereas conventional slicing retained across all $d$ and emitted $d+6$ nodes ($10$ vs.\ $6$ at $d=0,k=8$).

\paragraph{End-to-end custody and contract verification}
An integrated mock 322-slot dry run (316 provenance exclusions, six WSL2-executed family slots, 54 builds, 972 oracle runs, 90 profiles, 966 predictions) verified blinded score joining against an under-sized cohort. Complementary unit and live GPU checks in \nolinkurl{test_predict_external_contract.py} verify PEFT manifest hashing, all seven disposition statuses (\nolinkurl{RANKED}, \nolinkurl{NO_SYMPTOM}, \nolinkurl{EMPTY_SLICE}, \nolinkurl{NODE_LIMIT}, \nolinkurl{CONTEXT_LIMIT}, \nolinkurl{UNSUPPORTED_NODE_TOKEN}, \nolinkurl{NORMALIZATION_FAILURE}), and smoke 4-bit inference through pinned fold-0 adapter \texttt{33dfb098} against authored profiles into \nolinkurl{score_external.py} ($D_{\mathrm{pos}}=2$, $D_{\mathrm{elig}}=1$, 1/2 verified-negative alerts, recorded for the single completed live case prior to an earlier local inference service reset). Table~\ref{tab:dev-feasibility} consolidates all domain and boundary checks.

\subsection{Multi-agent prospective scalar study and blinded contract}
\label{sec:prospective-multiagent}

We executed bounded feasibility realizations of both prospective protocols across unmodified Linux WSL2 (Ubuntu, GCC~15.2.0, Bubblewrap~0.11.1, \nolinkurl{systemd-run}, FD~199 counter transport, and all four Gate~1 probes passing) and Windows~11 (MSYS2 GCC~15.1.0, RTX~4060 8\,GB GPU). Neither part reaches its confirmatory sample gate (Part~A \nolinkurl{primary_sized=false} at $10 < 60$ families and \nolinkurl{away_from_symptom_evidence=false} at $10 < 30$; Part~B \nolinkurl{positive_transfer_evidence=false} at $11 < 20$ groups), so all contrasts remain descriptive. Part~A produced byte-identical \nolinkurl{score.json} (\texttt{ce049b32}), \nolinkurl{frame.json}, \nolinkurl{draw.json}, and \nolinkurl{quarantine.json} across OSes, while \nolinkurl{labels.json}, \nolinkurl{blind.json}, and \nolinkurl{predictions.json} matched on all 322 dispositions, 42 projections, and 294 predictions (differing only on OS-specific \nolinkurl{operation_evidence_sha256} links). Four role-separated isolated subagents conducted Part~A verification and Part~B reference resolution, hash-chained into \nolinkurl{prospective_multiagent_results.json}.

In Part~A, an independent verifier checked 10 unseen C++17 scalar families (\texttt{F01}--\texttt{F10}, five linear-sham and five quadratic-sham, 42 active slots including the \texttt{F01} comparator swap) across plain, UBSan-trap, and counting builds with zero template collisions against the cohort or 77 pilot templates. Quarantine excluded 70 placeholder families (\texttt{E01}--\texttt{E70}, 280 slots bridged via \nolinkurl{exposure_frame_placeholders} among 102 exposure entries), while all 101 pilot and PIE exposure items had zero template or prompt-token collisions (21 live prompts vs.\ 38 \nolinkurl{abstract_graph} and 80 \nolinkurl{hotspot_raw} pilot prompts). Across 378 builds, 6,804 oracle runs, and 738 profiles over seeds 42--44, \texttt{work\_scoped} and \texttt{work\_conventional} core scores matched across seeds, while \texttt{hit\_growth} Top-1 was 10/20 on seed 42 and 9/20 on seeds 43--44 (Table~\ref{tab:prospective-multiagent}). On the 20 core positives (\nolinkurl{def_use=False}), \texttt{work\_scoped} scored 15/20 (75.00\%) Top-1 with 16/20 retention and 89.38\% reduction (matching \nolinkurl{hits_scoped} at 15/20 and raw order at 6/9 multi-candidate), compared with \texttt{work\_conventional} (0/20 Top-1, 20/20 retention, 64.34\% reduction; $+75.00$\,pp $\GroupedPartACoreCI$, LOO range $\GroupedPartACoreRange$) and pilot-refined \texttt{work\_scoped\_symbolic} (\nolinkurl{def_use=True}, 14/20, 70.00\%, inverted on \texttt{F08} off-site). The 0/20 CFG Top-1 stems from \nolinkurl{rank_candidates_by_role_priority} hoisting CFG-reachable \texttt{STATE\_UPD} initializers and accumulators ahead of the symptom on all 20 cases (including 10/10 at-site). Off-site edits scored 6/10 vs.\ 0/10 ($+60.00$\,pp $\GroupedPartAOffSiteCI$, missing \texttt{F02}, \texttt{F04}, \texttt{F06}, and two-hop \texttt{F10}); neither same-growth shams nor appropriate superlinear controls should alert. At the nominal threshold, false alerts occurred on \GroupedQuadraticFalse{} of quadratic shams and \GroupedLinearFalse{} of linear shams (\GroupedSameGrowthFalse{} overall), \GroupedAppropriateFalse{} of appropriate controls, and \GroupedSwapFalse{} of comparator-swap shams. Correct rejections on appropriate controls were \GroupedAppropriateCorrect{}; the replay receipt gives Clopper--Pearson intervals. On live 4-bit GPU inference (1.12--2.182\,GiB VRAM), fold-0 \nolinkurl{abstract_graph} scored 11/20 (55.00\%) constrained vs.\ Base 12/20 (60.00\%) and fold-0 \nolinkurl{hotspot_raw} 13/20 (65.00\%), a descriptive two-case difference on 10 families (with duplicate \texttt{F01} \nolinkurl{comparator_swap_positive}: 12/21, 11/21, and 14/21).

In Part~B, LF normalization aligned Windows GCC~15.1.0 (\texttt{2c95df0e}) with archived Linux GCC~15.2.0 (\texttt{c4260272}) across all 72 stress runs, 12 symptoms, slices, and rankings, and matched normalized hashes on \texttt{PIE\_DEV\_p03497} and \texttt{AUTHORED\_SMOKE\_POS} (with \nolinkurl{quarantine_sha256} hashing the fixed tag \nolinkurl{part-b-multiagent-quarantine-11-groups}). The cohort comprises 12 pilot pairs (including four \texttt{G\_STRESS\_08} pairs) held strictly out-of-fold, one authored cursor-reset pair, \texttt{PIE\_DEV\_p03497} (timing slope $<0.50$ on all three seeds), and two verified negatives ($14+2$ pairs in $11+2$ groups). Reviewers~A and~B marked identical spans on 13/14 positives (\texttt{PAIR\_013} by consensus, 13/14 via Adjudicator~C). Across $D_{\mathrm{pos}}=14$ ($D_{\mathrm{elig}}=14$), \texttt{work\_scoped} retained 12/14 (81.82\% group; 86.47\% reduction, missing \texttt{PAIR\_039} and \texttt{PAIR\_089}). Strict OOF \nolinkurl{abstract_graph} scored 10/14 (63.64\% group; vs.\ 12/14 = 81.82\% in-fold \texttt{fold-0}), compared with unbiased \texttt{symbolic} (\nolinkurl{def_use=False}, 6/14, 54.55\%; $+9.09$\,pp $\GroupedPartBUnbiasedCI$, 6 vs.\ 2 pairs) and matching refined \nolinkurl{prediction_contract.py} \texttt{symbolic} (\nolinkurl{def_use=True}, 10/14, 63.64\%; 0.00\,pp $\GroupedPartBRefinedCI$). All three methods hit both non-pilot positives, so differences come from the 12 held-out pilot pairs; excluding sub-threshold \texttt{PIE\_DEV\_p03497} gives 9/13 for strict OOF \nolinkurl{abstract_graph} and refined \texttt{symbolic} vs.\ 5/13 for unbiased \texttt{symbolic}. On $D_{\mathrm{rank}}=10$, strict OOF \nolinkurl{abstract_graph} and refined \texttt{symbolic} scored 8/10 (71.43\%) vs.\ 4/10 (57.14\%) for unbiased \nolinkurl{def_use=False}, with 1/2 negative alerts.

\begin{table}[t]
\centering\scriptsize
\setlength{\tabcolsep}{2pt}
\caption{Bounded multi-agent prospective scalar cohort (Part~A, \texttt{primary\_sized=false}) and blinded adjudication contract (Part~B, \texttt{positive\_transfer\_evidence=false}). Primary metric within these secondary analyses: equal-group Top-1. GPU runs use 4-bit NF4 \texttt{Qwen2.5-Coder-1.5B} (max $2.182$\,GiB VRAM).}
\label{tab:prospective-multiagent}
\begin{tabular}{@{}llrr@{}}
\toprule
Cohort & Method / Condition & Ret. & Top-1 (Grp \%) \\
\midrule
A (20 core) & \texttt{work\_scoped} (\texttt{def\_use=F}) & 16/20 & 15/20 (75.00) \\
A (20 core) & \texttt{work\_conventional} (\texttt{def\_use=F}) & 20/20 & 0/20 (0.00) \\
A (20 core) & \texttt{work\_scoped\_symbolic} (\texttt{def\_use=T}) & 16/20 & 14/20 (70.00) \\
A (20 core) & Base (\texttt{abstract\_graph}, Con.) & 16/20 & 12/20 (60.00) \\
A (20 core) & QLoRA (\texttt{abstract\_graph}, Unc.) & 20/20 & 2/20 (10.00) \\
A (20 core) & QLoRA (\texttt{abstract\_graph}, Con.) & 16/20 & 11/20 (55.00) \\
A (20 core) & QLoRA (\texttt{hotspot\_raw}, Con.) & 16/20 & 13/20 (65.00) \\
\midrule
B ($D_{\mathrm{pos}}$) & Unbiased \texttt{symbolic} (\texttt{def\_use=F}) & 12/14 & 6/14 (54.55) \\
B ($D_{\mathrm{pos}}$) & Refined \texttt{symbolic} (\texttt{def\_use=T}) & 12/14 & 10/14 (63.64) \\
B ($D_{\mathrm{pos}}$) & Fold-0 \texttt{abstract\_graph} (in-fold) & 12/14 & 12/14 (81.82) \\
B ($D_{\mathrm{pos}}$) & Strict OOF \texttt{abstract\_graph} & 12/14 & 10/14 (63.64) \\
B ($D_{\mathrm{rank}}$) & Unbiased \texttt{symbolic} (\texttt{def\_use=F}) & 10/10 & 4/10 (57.14) \\
B ($D_{\mathrm{rank}}$) & Refined \texttt{symbolic} (\texttt{def\_use=T}) & 10/10 & 8/10 (71.43) \\
B ($D_{\mathrm{rank}}$) & Strict OOF \texttt{abstract\_graph} & 10/10 & 8/10 (71.43) \\
\bottomrule
\end{tabular}
\end{table}
